\documentclass[a4paper,12pt]{article}
\usepackage[T2A]{fontenc}
\usepackage[utf8]{inputenc}
\usepackage[russian]{babel}
\usepackage{amsmath,amssymb,amsfonts}
\usepackage{graphicx}
\usepackage{float}
\usepackage{caption}
\usepackage{subcaption}
\usepackage{geometry}
\geometry{left=3cm, right=1.5cm, top=2cm, bottom=2cm}
\usepackage{setspace}
\onehalfspacing % Полуторный интервал
\usepackage{indentfirst} % Красная строка
\setlength{\parindent}{1.25cm} % Отступ первой строки
\usepackage{enumitem}
\usepackage{booktabs} % Для красивых таблиц
\usepackage{algorithm}
\usepackage{algpseudocode}
\usepackage{hyperref}
\usepackage{tikz}
\usetikzlibrary{trees, positioning}
\hypersetup{
    colorlinks=true,
    linkcolor=blue,
    filecolor=magenta,
    citecolor=red,     
    urlcolor=cyan,
    pdftitle={Курсовой проект},
    pdfpagemode=FullScreen,
}


\usepackage{amsthm} % Для теорем, лемм, определений
\usepackage{thmtools} % Для дополнительного оформления
\usepackage{xcolor} % Для цветов

% Математические обозначения
\newcommand{\R}{\mathbb{R}}
\newcommand{\norm}[1]{\left\lVert #1 \right\rVert}
\newcommand{\fro}[1]{\left\lVert #1 \right\rVert_{\mathrm F}}
\newcommand{\nuc}[1]{\left\lVert #1 \right\rVert_*}
\newcommand{\tr}{\operatorname{tr}}
\newcommand{\rank}{\operatorname{rank}}
\newcommand{\grad}{\operatorname{grad}}
\newcommand{\argmin}{\operatorname*{arg\,min}}
\renewcommand{\Proj}{P}
\newcommand{\Mcal}{\mathcal{M}}
\newcommand{\Retr}{\operatorname{Retr}}
\newcommand{\St}{\mathrm{St}}
\newcommand{\Gr}{\mathrm{Gr}}

% Определяем стили для окружений
\declaretheoremstyle[
    spaceabove=6pt, spacebelow=6pt,
    headfont=\normalfont\bfseries,
    notefont=\mdseries, notebraces={(}{)},
    bodyfont=\normalfont,
    postheadspace=1em,
    headpunct={}
]{mystyle}

% Создаем окружения с нумерацией внутри разделов
\declaretheorem[style=mystyle, name=Определение, numberwithin=section]{definition}
\declaretheorem[style=mystyle, name=Теорема, sibling=definition]{theorem}
\declaretheorem[style=mystyle, name=Лемма, sibling=definition]{lemma}
\declaretheorem[style=mystyle, name=Следствие, sibling=definition]{corollary}
\declaretheorem[style=mystyle, name=Утверждение, sibling=definition]{proposition}
\declaretheorem[style=mystyle, name=Пример, numbered=no]{example}



\begin{document}

% Титульный лист
\begin{titlepage}
    \centering
    {\large\bfseries Национальный исследовательский университет \\ «Высшая школа экономики»\par}
    \vspace{1cm}
    {\large Факультет экономических наук, Факультет компьютерных наук\par}
    {\large Образовательные программы «Экономика и анализ данных», «Прикладная математика и информатика»\par}
    \vspace{2cm}
    {\Large\bfseries КУРСОВОЙ ПРОЕКТ\par}
    \vspace{1cm}
    {\Large\bfseries Геометрические методы на многообразии матриц фиксированного ранга для заполнения матриц \par}
    \vspace{2cm}
    \begin{flushright}
        \large
        \textbf{Выполнили:}\\
        Студент 2 курса, группы БЭАД245\\
        Аляутдинов Карим Маратович,\\
        Студент 2 курса, группы БПМИИ245\\
        Исмаилов Магомед Айдынович.\\
        \vspace{0.5cm}
        \textbf{Научный руководитель:}\\
        профессор, д.ф.-м.н.\\
        Широков Дмитрий Сергеевич\\
    \end{flushright}
    \vfill
    {\large Москва, 2026\par}
\end{titlepage}

% Содержание
\tableofcontents
\newpage

% Основной текст
\section*{Введение}
\addcontentsline{toc}{section}{Введение}
\phantomsection

\textbf{Актуальность темы.} \\[5pt] Тема восстановления матриц наиболее актуальна в современных прикладных задачах, таких как рекомендательные системы, обработка изображений, восстановление панельных экономических и социальных данных. \\[5pt] Зачастую данные в восстанавливаемых матрицах играют ключевую роль исследований, потому что на них строится дальнейший анализ. На качество восстановленных данных влияет устойчивость метода к шуму, количество пропусков и их структура в матрице, истинный ранг и множество других факторов. \\[5pt] Также во многих прикладных задачах, в том числе в области экономики (например: макроэкономический мониторинг, алготрейдинг, кредитный скоринг, динамическое ценообразование и т.д.) используются low-latency алгоритмы, поэтому большое значение имеет скорость сходимости метода восстановления. \\[5pt] Существенная часть проблемы состоит в том, что традиционные методы восстановления больших матриц требуют дорогостоящих операций, например convex-relaxation использует SVD на каждой итерации, а ALS считает МНК и QR разложение. Они довольно тяжелые с вычислительной точки зрения. \\[5pt] На этом фоне актуальным становится исследование методов, которые используют предположение о ранге матрицы для улучшеной скорости сходимости и устойчивости к шуму. В нашем случае одним из таких методов будет риманов градиентный спуск на многообразии матриц фиксированного ранга.

\textbf{Цель исследования.} \\[5pt] Реализовать и протестировать метод риманова градиентного спуска для завершения матриц. Сравнить его эффективность и скорость с классическим методом минимизации ядерной нормы и геометрическими подходами (ALS, OptSpace) на реальных экономических данных.

\textbf{Задачи исследования.} \\[5pt]
1. Провести обзор литературы по методам заполнения матриц. \\[5pt]
2. Описать геометрическую постановку задачи на многообразии
матриц фиксированного ранга и вывести шаг риманова градиентного спуска. \\[5pt]
3. Реализовать алгоритм риманова градиентного спуска. \\[5pt]
4. Воспроизвести методы минимизации ядерной нормы, ALS и OptSpace. \\[5pt]
5. Подготовить экономические данные для экспериментов. \\[5pt]
6. Провести сравнительные эксперименты по качеству восстановления и вычислительной эффективности. \\[25pt]
7. Исследовать устойчивость методов к шуму, доле и структуре пропусков, влияние неверного ранга и чувствительность к
инициализации. \\[5pt]
8. Проанализировать полученные результаты и выявить режимы преимуществ и недостатков риманова градиентного спуска в задаче восстановления матриц. \\[5pt]
\textbf{Объект и предмет исследования.}
\begin{itemize}
    \item \textbf{Объект:} процессы восстановления пропущенных значений в низкоранговых матрицах экономических данных.
    \item \textbf{Предмет:} методы оптимизации на многообразии матриц фиксированного ранга (в частности, риманов градиентный спуск) и их свойства: точность восстановления, скорость сходимости и устойчивость к шуму/инициализации по сравнению с классическими подходами.
\end{itemize}
\textbf{Методы исследования.} \\[5pt]
1. Анализ литературы по matrix completion и римановой оптимизации. \\[5pt]
2. Численные методы оптимизации. \\[5pt]
3. Вычислительный эксперимент на реальных экономических данных. \\[5pt]
4. Методы статистической оценки качества восстановления: $RMSE$, $MAE$, $||\cdot||_F.$ \\[5pt]
5. Анализ вычислительной эффективности: время выполнения, число итераций. \\[5pt]
6. Анализ устойчивости (sensitivity analysis) к шуму, доле/структуре пропусков, инициализации и выбору ранга. \\[5pt]
\textbf{Структура работы.} \\[5pt]
Структура работы включает введение, теоретическую и практическую части, заключение, список литературы и приложения. Во введении обосновывается актуальность темы, формулируются цель, задачи, объект и предмет исследования, а также описываются применяемые методы. \\[5pt]
В первом разделе представлен обзор литературы по задаче заполнения матриц: рассматриваются методы выпуклой релаксации, алгоритмы факторизации и геометрические подходы на многообразиях, а также их преимущества и ограничения. Во втором разделе вводятся основные определения, обозначения и математический аппарат, используемый в работе. \\[5pt]
Третий раздел посвящен теоретической постановке задачи matrix completion на многообразии матриц фиксированного ранга и описанию метода риманова градиентного спуска. В четвертом разделе описывается программная реализация, данные, дизайн вычислительных экспериментов и результаты сравнительного анализа с baseline-методами (минимизация ядерной нормы, ALS, OptSpace) по качеству восстановления, скорости сходимости и устойчивости к шуму и инициализации. \\[25pt]
В заключении формулируются основные выводы, ограничения исследования и направления дальнейшей работы. В приложениях приводятся дополнительные таблицы, графики и технические детали реализации.

\newpage

\section{Обзор литературы}
\label{sec:literature}

\subsection{Классическая выпуклая линия}

Работы Candes и Recht, Candes и Tao, Recht формируют теоретический фундамент выпуклой релаксации через ядерную норму. В этих работах они от постанвоки задачи восстановления матриц как выпуклой оптимизации приходят к почти-оптимальным условиям (с точки зрения числа наблюдаемых элементов) решения задачи. В нашей работе выпуклая постановка рассматривается как необходимая теория для алгоритмов SVT, SoftImpute.

\subsection{Алгоритмы для больших матриц}

В статьях Cai--Cand\`es--Shen и Mazumder--Hastie--Tibshirani теоретия релаксации ядерной нормы прикладывается в восстановлении матриц через разработанные алгоритмы SVT и SoftImpute, с которыми ведется основное сравнение в нашей работе. Их преимущество $-$ устойчивое поведение и хорошая инженерная применимость на разреженных матрицах. Для нас это ключевой сравнительный пункт: если геометрический метод не дает выигрыша по скорости или устойчивости, его использование теряет смысл.

\subsection{Шум и более реалистичная постановка}

Работа Candes и Plan направлена на то, чтобы сформулировать ппроблему в реальных условиях, так как обычно восстанавливаемые данные сопровождаются шумом, это делает сравнение методов более прикладным. Одновременно остается нерешенный вопрос: как именно меняется поведение разных алгоритмов при одинаковом уровне шума и при разных долях наблюдений.

\subsection{Невыпуклые и фиксированно-ранговые методы}

Статья Keshavan--Montanari--Oh является центральной в нашей теме. В ней авторы впервые переходят от релаксации ядерной нормы к более геометрическому подходу. Разработанный алгоритм OptSpace использует предположение о ранге для нахождения матрицы в многообразии матриц фиксированного ранга. Преимущество подхода заключается в вычислительной экономии. Ограничения - более высокая чувствительность к инициализации и параметрам. Для нашей работы это не недостаток, а предмет анализа: именно в таких условиях интересно проверять, насколько геометрические методы уменьшают нестабильность невыпуклой оптимизации.

\subsection{Расширения и прикладной контекст}

Работы Jain--Dhillonи Athey et al. показывают, что практические задачи часто выходят за рамки ``чистой'' матричной постановки: появляются внешние признаки, панельная структура и прикладные ограничения интерпретации. Обзор Davenport--Romberg фиксирует общий вывод: универсально лучшего метода нет, а выбор всегда задается компромиссом между гарантиями, ценой вычислений и гибкостью модели. Это и определяет дизайн нашей работы: мы сравниваем методы в едином экспериментальном протоколе, а не пытаемся заранее объявить один подход доминирующим.

\subsection{Выводы из обзора литературы}

Из обзора литературы следует, что выпуклые методы важны как теоретически обоснованный ориентир, но в больших задачах их вычислительная стоимость часто оказывается ограничивающим фактором. Невыпуклые и геометрические подходы, напротив, лучше масштабируются, однако сильнее зависят от инициализации, выбора ранга и параметров алгоритма. Поэтому в данной работе риманов градиентный спуск рассматривается как основной метод, а его эффективность оценивается в сравнении с базовыми подходами по качеству восстановления, скорости сходимости и устойчивости к шуму.


\begin{table}[H]
    \centering
    \caption{Сводная таблица исследований по теме}
    \label{tab:literature}
    \begin{tabular}{p{0.15\linewidth}p{0.25\linewidth}p{0.2\linewidth}p{0.3\linewidth}}
        \toprule
        \textbf{Автор(ы), год}                                                                                       & \textbf{Цель/Гипотеза}                                                     & \textbf{Данные и методы}                                            & \textbf{Ключевые результаты}                                                                                                                \\
        \midrule
        \href{https://arxiv.org/abs/0805.4471}{Cand\`es, Recht (2008)}                                               & Обосновать точное восстановление низкоранговой матрицы по части наблюдений & Выпуклая релаксация ранга через ядерную норму; случайные наблюдения & Гарантии восстановления при условиях инкогерентности и случайных пропусках.                                                                 \\
        \href{https://arxiv.org/abs/0805.4471}{Cand\`es, Tao (2009)}                                                 & Уточнить условия восстановления и оценить требуемое число наблюдений       & Теоретический анализ выпуклой постановки                            & Близкие к оптимальным оценки объема наблюдений и условие сильной инкогерентности для точного восстановления.                                \\
        \href{https://www.jmlr.org/papers/volume12/recht11a/recht11a.pdf}{Recht (2011)}                              & Упростить доказательства гарантий восстановления, ослабить условия         & Теоретический анализ выпуклой модели                                & Более компактное доказательство результатов, близких к работам Cand\`es--Tao, доказательство достаточности условия обычной инкогерентности. \\
        \href{https://arxiv.org/abs/0810.3286}{Cai, Cand\`es, Shen (2008)}                                           & Сделать выпуклую модель вычислительно применимой на больших матрицах       & Алгоритм Singular Value Thresholding (SVT)                          & Алгоритм решения задачи matrix completion.                                                                                                  \\
        \href{https://www.jmlr.org/papers/volume11/mazumder10a/mazumder10a.pdf}{Mazumder, Hastie, Tibshirani (2010)} & Разработать масштабируемый метод для больших и разреженных матриц          & Спектральная регуляризация, алгоритм Soft-Impute                    & Практическая эффективность на больших разреженных данных.                                                                                   \\
        \href{https://arxiv.org/abs/0903.3131}{Cand\`es, Plan (2009)}                                                & Перейти к постановке с шумом                                               & Теоретический анализ matrix completion с шумом                      & Результаты по устойчивости восстановления при шумных наблюдениях.                                                                           \\
        \href{https://arxiv.org/abs/0901.3150}{Keshavan, Montanari, Oh (2009)}                                       & Исследовать невыпуклое восстановление с фиксированным рангом               & Факторизационный метод OptSpace                                     & Алгоритм OptSpace фиксированно-ранговой постановки.                                                                                         \\
        \href{https://arxiv.org/abs/1306.0626}{Jain, Dhillon (2013)}                                                 & Включить дополнительную информацию о строках и столбцах                    & Inductive matrix completion с признаками объектов                   & Показано, что признаки улучшают восстановление.                                                                                             \\
        \href{https://arxiv.org/abs/1601.06422}{Davenport, Romberg (2016)}                                           & Систематизировать подходы low-rank recovery                                & Обзор литературы по восстановлению неполных наблюдений              & Показан компромисс между теоретическими гарантиями, вычислительной стоимостью и гибкостью моделей.                                          \\
        \bottomrule
    \end{tabular}
\end{table}
\begin{table}[H]
    \centering
    \caption{Сводная таблица исследований по теме}
    \label{tab:literature}
    \begin{tabular}{p{0.15\linewidth}p{0.25\linewidth}p{0.2\linewidth}p{0.3\linewidth}}
        \toprule
        \textbf{Автор(ы), год}                                       & \textbf{Цель/Гипотеза}                                                & \textbf{Данные и методы}                                         & \textbf{Ключевые результаты}                                                                          \\
        \midrule
        \href{https://arxiv.org/abs/1710.10251}{Athey et al. (2017)} & Применить matrix completion в причинном анализе панельных данных      & Методы matrix completion для causal panel data models            & Продемонстрирована прикладная ценность метода в эконометрике и анализе панелей.                       \\
        \href{https://arxiv.org/abs/1910.06677}{Bai, Ng (2019)}      & Связать matrix completion, контрфактуалы и факторный анализ пропусков & Панельные модели с пропусками; контрфактуальный анализ & Показана связь между matrix completion и факторным подходом для работы с отсутствующими наблюдениями. \\
        \bottomrule
    \end{tabular}
\end{table}

\subsection{Граф развития темы восстановления матриц}
\vspace{1cm}
\begin{tikzpicture}[
        level distance=2cm,
        sibling distance=3.5cm,
        every node/.style={draw=none, align=center}
    ]

    \node {Convex Theory\\[5pt]Cand\`es--Recht}
    child {
            node {Cand\`es--Tao}
        }
    child {
            node {Soft-Impute}
            child{
                    node{SVT}
                }
            child{
                    node{Cand\`es--Plan}
                }
        }
    child {
            node {Recht}
        };

    \node (opt) [below=3cm of current bounding box.south] {Fixed rank\\+ Grassmann geometry\\[5pt]Opt Space};

    \draw (current bounding box.south) -- (opt);

    \node (vd) [below left=of opt] {Jain--Dhillon\\(side info)};
    \node (ai) [below right=of opt] {Ashley et al.\\(causal panel data)};

    \draw (opt) -- (vd);
    \draw (opt) -- (ai);

    \node (riem) [below=of opt, yshift=-2cm] {Riemannian\\fixed-rank GD on manifold};

    \draw (opt) -- (riem);

\end{tikzpicture}

\newpage

\section{Основные определения и обозначения}
\label{sec:definitions}

Пусть дана матрица
\[
    Y \in \R^{m \times n},
\]
в которой наблюдаются только элементы с индексами из множества
\[
    \Omega \subseteq \{1,\dots,m\} \times \{1,\dots,n\}.
\]
Предполагается, что существует скрытая матрица
\[
    X_* \in \R^{m \times n},
\]
имеющая малый ранг, и наблюдения удовлетворяют соотношению
\[
    Y_{ij} = (X_*)_{ij} + \varepsilon_{ij}, \qquad (i,j)\in\Omega,
\]
где $\varepsilon_{ij}$ описывает шум.

\begin{definition}
    Введем маску наблюдений $W\in\R^{m\times n}$:
    \[
        W_{ij} =
        \begin{cases}
            1, & (i,j)\in\Omega,    \\
            0, & (i,j)\notin\Omega.
        \end{cases}
    \]
    Тогда оператор наблюдаемых элементов удобно записывать через адамарово произведение:
    \[
        \Proj_\Omega(X)=W\odot X.
    \]
    Здесь $\odot$ означает покомпонентное произведение матриц.
\end{definition}

Тогда задача matrix completion состоит в восстановлении матрицы $X_*$ по наблюдаемой части $\Proj_\Omega(Y)$ при предположении, что искомая матрица имеет малый ранг.

В дальнейшем используются следующие обозначения:
\begin{itemize}
    \item $\fro{X}$ --- норма Фробениуса матрицы $X$;
    \item $\nuc{X}$ --- ядерная норма матрицы $X$;
    \item $\rank(X)$ --- ранг матрицы $X$;
    \item $\Mcal_r = \{X \in \R^{m\times n} : \rank(X)=r\}$ --- множество матриц фиксированного ранга $r$;
    \item $U\Sigma V^\top$ --- компактное сингулярное разложение матрицы ранга $r$.
\end{itemize}

\newpage

\section{Теоретическая часть}
\label{sec:theory}

\subsection{Карты, атласы и гладкость}

В этом разделе используется только тот фрагмент дифференциальной геометрии, который нужен для постановки риманова градиентного спуска. Материал опирается на лекционные и семинарские курсы по дифференциальной геометрии А. О. Иванова, А. Т. Фоменко и А. В. Пенского \cite{ivanov-dg,fomenko-dg,penskoy-seminars}.

\begin{definition}
    Пусть $M$ --- топологическое пространство. Картой размерности $d$ называется пара $(U,\varphi)$, где $U\subset M$ --- открытое множество, а
    \[
        \varphi:U\to \varphi(U)\subset\R^d
    \]
    --- гомеоморфизм на открытое множество $\varphi(U)$.
\end{definition}

Карта позволяет заменить точки из $U$ их координатами в $\R^d$. Если таких карт достаточно, чтобы покрыть все множество, мы получаем атлас.

\begin{definition}
    Набор карт $\mathcal A=\{(U_\alpha,\varphi_\alpha)\}$ называется атласом, если
    \[
        M=\bigcup_\alpha U_\alpha.
    \]
    Если все карты имеют значения в $\R^d$, то число $d$ называется размерностью многообразия.
\end{definition}

\begin{definition}
    Топологическое пространство $M$ называется $d$-мерным топологическим многообразием, если у каждой точки $p\in M$ есть окрестность, гомеоморфная открытому множеству в $\R^d$.
\end{definition}

Иначе говоря, многообразие может быть устроено глобально сложно, но локально оно выглядит как обычное евклидово пространство. Именно поэтому на нем можно вводить локальные координаты.

Если две карты пересекаются, то одна и та же точка получает две записи. Поэтому важно отображение перехода между координатами:
\[
    \psi\circ\varphi^{-1}:\varphi(U\cap V)\to\psi(U\cap V).
\]

\begin{definition}
    Атлас называется гладким, если все отображения перехода между его картами гладкие. Множество с таким атласом называется гладким многообразием.
\end{definition}

\begin{definition}
    Пусть $M$ и $N$ --- гладкие многообразия. Отображение $F:M\to N$ называется гладким, если в любых картах его координатная запись
    \[
        \psi\circ F\circ\varphi^{-1}
    \]
    является гладким отображением между открытыми подмножествами евклидовых пространств.
\end{definition}

Гладкость нам нужна для того чтобы производные, касательные векторы и градиенты не зависели от произвольного выбора координат.

\subsection{Вложенные многообразия и регулярные уровни}

В оптимизации чаще всего встречаются не абстрактные многообразия, а подмножества евклидова пространства, заданные ограничениями. Такие множества называют вложенными многообразиями.

\begin{definition}
    Подмножество $M\subset\R^N$ называется вложенным гладким многообразием размерности $d$, если в окрестности каждой своей точки оно после гладкой замены координат выглядит как $\R^d\times\{0\}^{N-d}$.
\end{definition}

Один из стандартных способов получить вложенное многообразие --- задать его системой уравнений. Для этого нужны регулярные точки.

\begin{definition}
    Пусть $F:\R^N\to\R^k$ --- гладкое отображение. Точка $x$ называется регулярной точкой, если матрица Якоби $DF(x)$ имеет полный ранг $k$. Если ранг меньше $k$, точка называется особой.
\end{definition}

\begin{definition}
    Точка $y\in\R^k$ называется регулярным значением отображения $F$, если каждая точка из прообраза $F^{-1}(y)$ является регулярной. Если в прообразе есть хотя бы одна особая точка, то $y$ называется особым значением.
\end{definition}

\begin{theorem}[Теорема о регулярном уровне]
    Пусть $F:\R^N\to\R^k$ --- гладкое отображение, и для всех $x$ таких, что $F(x)=0$, матрица $DF(x)$ имеет ранг $k$. Тогда множество
    \[
        M=\{x\in\R^N:F(x)=0\}
    \]
    является вложенным гладким многообразием размерности $N-k$.
\end{theorem}

Это утверждение является прямым следствием теоремы о неявной функции: если уравнения локально независимы, то часть координат можно выразить через остальные. Поэтому остается $N-k$ свободных координат.

Обратная идея тоже важна: локально вложенное многообразие можно описывать системой гладких уравнений с полным рангом Якоби. Для нашей темы это объясняет, почему ограничения в оптимизации можно изучать геометрически, а не только как набор формальных равенств.

\subsection{Касательные векторы}

Касательный вектор сначала можно определить через координаты. Пусть $x\in M$ и выбрана карта $(U,\varphi)$ около $x$. Тогда вектор в координатах --- это элемент $v\in\R^d$. Если перейти к другой карте $(V,\psi)$, тот же вектор должен измениться по правилу
\[
    v_\psi =
    D(\psi\circ\varphi^{-1})(\varphi(x))\,v_\varphi.
\]
Поэтому определение касательного вектора не зависит от выбора конкретной карты: при смене координат его координатное представление меняется по правилу дифференциала отображения перехода.

Для наших целей удобнее эквивалентное определение через кривые:
\[
    T_xM=\{\dot\gamma(0):\gamma(t)\in M,\ \gamma(0)=x\}.
\]
Существует стандартная теорема об эквивалентности координатного определения и определения через кривые; ее доказательство здесь не приводится. В дальнейшем мы будем пользоваться определением через кривые, потому что оно удобнее для вложенных многообразий.

\subsection{Риманова метрика и градиент}

\begin{definition}
    Риманова метрика --- это выбор скалярного произведения $\langle\cdot,\cdot\rangle_x$ на каждом касательном пространстве $T_xM$, гладко зависящий от точки $x$. Для любых $\xi,\eta,\zeta\in T_xM$ и чисел $a,b$ оно удовлетворяет:
    \[
        \langle a\xi+b\eta,\zeta\rangle_x
        =a\langle\xi,\zeta\rangle_x+b\langle\eta,\zeta\rangle_x,
    \]
    \[
        \langle \xi,\eta\rangle_x=\langle\eta,\xi\rangle_x,\qquad
        \langle\xi,\xi\rangle_x>0\quad \text{при }\xi\ne 0.
    \]
\end{definition}

Метрика нужна для того, чтобы измерять длины и углы касательных векторов. Без нее нельзя однозначно сказать, какое касательное направление является направлением наибольшего возрастания функции.

В этой работе используется метрика, индуцированная нормой Фробениуса. Для матриц одинакового размера:
\[
    \langle A,B\rangle=\tr(A^\top B).
\]
Соответствующая норма:
\[
    \fro{A}=\sqrt{\langle A,A\rangle}.
\]
Такой выбор естественен, потому что сама функция потерь в задаче matrix completion записывается через норму Фробениуса.

Пусть $f:M\to\R$ --- гладкая функция. Ее дифференциал $Df(x)$ является линейной функцией на $T_xM$: он принимает касательное направление $\xi$ и возвращает производную функции вдоль этого направления.

\begin{definition}
    Риманов градиент $\grad f(x)$ определяется условием
    \[
        Df(x)[\xi]=\langle \grad f(x),\xi\rangle_x
        \qquad \text{для всех }\xi\in T_xM.
    \]
\end{definition}

Таким образом, градиент определяется не только функцией, но и выбранной метрикой: именно метрика отождествляет линейную форму $Df(x)$ с касательным вектором.

Если $M$ вложено в евклидово пространство и метрика индуцирована обычным скалярным произведением, то риманов градиент получается проекцией евклидова градиента на касательное пространство:
\[
    \grad_M f(x)=\Proj_{T_xM}(\nabla f(x)).
\]
Нормальная компонента евклидова градиента не задает допустимое движение вдоль многообразия. Поэтому в алгоритме остается только касательная часть. Именно эта формула будет использоваться дальше для матриц фиксированного ранга.

\subsection{Геодезические и ретракции}

На многообразии градиент $\grad f(x)$ лежит в касательном пространстве $T_xM$, а не в самом многообразии. Поэтому выражение
\[
    x-\eta \grad f(x)
\]
обычно не является точкой из $M$. Нужно сначала сделать шаг в касательном направлении, а затем вернуться на допустимое множество.

Самый геометрически точный способ --- двигаться по геодезической. Геодезическая играет роль прямой линии на многообразии: ее начальная точка равна $x$, а начальная скорость равна выбранному касательному вектору. Если обозначить через $\operatorname{Exp}_x$ экспоненциальное отображение, то идеальный шаг можно записать так:
\[
    x_{k+1}=\operatorname{Exp}_{x_k}\bigl(-\eta_k\grad f(x_k)\bigr).
\]
Проблема в том, что экспоненциальное отображение часто трудно вычислять явно. Поэтому в алгоритмах обычно используют более простую замену --- ретракцию.

\begin{definition}
    Ретракцией называется отображение $\Retr_x:T_xM\to M$, которое удовлетворяет двум условиям:
    \[
        \Retr_x(0_x)=x,
        \qquad
        D\Retr_x(0_x)=\operatorname{id}_{T_xM}.
    \]
\end{definition}

Первое условие говорит, что нулевой шаг оставляет точку на месте. Второе означает, что при малых шагах ретракция не искажает касательное направление в первом порядке. Поэтому ретракция может заменить геодезическую в градиентном методе: она сохраняет нужную локальную геометрию, но обычно считается проще.

Для матриц фиксированного ранга естественная ретракция будет устроена так: сделать касательный шаг, а затем заменить полученную матрицу ближайшей матрицей ранга $r$. Ниже это приведет к усеченному SVD.

\subsection{Постановка задачи}

В задаче заполнения матриц наблюдается только часть элементов матрицы $Y$. Классическая выпуклая постановка использует ядерную норму:
\begin{equation}
    \label{eq:convex-final}
    \min_{X \in \R^{m\times n}}
    \frac12 \fro{\Proj_\Omega(X-Y)}^2 + \lambda \nuc{X}.
\end{equation}
Первое слагаемое отвечает за согласование с наблюдаемыми элементами, второе поощряет малый ранг. Эта модель выпуклая, но в больших задачах может быть вычислительно дорогой.

В fixed-rank подходе ранг задается заранее, и оптимизация ведется напрямую по множеству матриц ранга $r$:
\begin{equation}
    \label{eq:fixed-rank-final}
    \min_{X \in \Mcal_r}
    \frac12 \fro{\Proj_\Omega(X-Y)}^2,
    \qquad
    \Mcal_r=\{X\in\R^{m\times n}:\rank(X)=r\}.
\end{equation}

Для шумных данных дополнительно рассматривается L2-регуляризация:
\begin{equation}
    \label{eq:fixed-rank-regularized}
    \min_{X \in \Mcal_r} f_\lambda(X)
    = \frac12 \fro{\Proj_\Omega(X-Y)}^2 + \frac{\lambda}{2}\fro{X}^2,
    \qquad \lambda\ge 0.
\end{equation}
Штраф $\fro{X}^2$ ограничивает размер решения и снижает риск подстройки под шум.

\subsection{Риманов градиентный спуск}

Рассмотрим функцию
\[
    f(X)=\frac12\fro{\Proj_\Omega(X-Y)}^2.
\]
Ее евклидов градиент равен
\[
    \nabla f(X)=\Proj_\Omega(X-Y).
\]
Для регуляризованной функции \eqref{eq:fixed-rank-regularized} получаем
\[
    \nabla f_\lambda(X)=\Proj_\Omega(X-Y)+\lambda X.
\]
Это следует из дифференцирования по произвольному приращению $H$:
\[
    \left.\frac{d}{dt}\right|_{t=0}
    \frac12\fro{\Proj_\Omega(X+tH-Y)}^2
    =
    \langle \Proj_\Omega(X-Y),\Proj_\Omega(H)\rangle.
\]
Так как $\Proj_\Omega$ является ортогональной проекцией,
\[
    \langle \Proj_\Omega(X-Y),\Proj_\Omega(H)\rangle
    =
    \langle \Proj_\Omega(X-Y),H\rangle.
\]
Отсюда получается формула для $\nabla f(X)$, а регуляризация добавляет градиент $\lambda X$.

Риманов градиент --- это проекция евклидова градиента на касательное пространство:
\begin{equation}
    \label{eq:riemannian-gradient-final}
    \grad_{\Mcal_r} f_\lambda(X)
    =
    \Proj_{T_X\Mcal_r}\bigl(\Proj_\Omega(X-Y)+\lambda X\bigr).
\end{equation}

После касательного шага нужно вернуться на $\Mcal_r$. Для этого используется ретракция:
\begin{equation}
    \label{eq:fixed-rank-retraction-final}
    \Retr_X(\Xi)=\Pi_r(X+\Xi),
\end{equation}
где $\Pi_r$ --- усеченное SVD. Если
\[
    X+\Xi=\sum_i \sigma_i u_i v_i^\top
\]
--- сингулярное разложение, то
\[
    \Pi_r(X+\Xi)=\sum_{i=1}^r \sigma_i u_i v_i^\top.
\]
По теореме Эккарта--Янга это ближайшая матрица ранга не выше $r$ в норме Фробениуса. В окрестности точки ранга $r$ такая операция возвращает нас на $\Mcal_r$.

Проверим условия ретракции. При $\Xi=0$ имеем $\Pi_r(X)=X$. Кроме того, для малого $\Xi\in T_X\Mcal_r$:
\[
    \Pi_r(X+\Xi)=X+\Xi+O(\fro{\Xi}^2).
\]
Следовательно, дифференциал ретракции в нуле является тождественным на $T_X\Mcal_r$. Геодезики здесь не используются: для градиентного метода достаточно совпадения с касательным шагом в первом порядке, а усеченное SVD считается проще.

Итоговый шаг риманова градиентного спуска:
\begin{equation}
    \label{eq:rgd-final}
    X_{k+1}
    =
    \Retr_{X_k}\bigl(-\eta_k\grad_{\Mcal_r} f_\lambda(X_k)\bigr).
\end{equation}
Длина шага $\eta_k$ подбирается фиксированно или с помощью backtracking line search.

\begin{algorithm}[H]
    \caption{Риманов градиентный спуск на $\Mcal_r$}
    \label{alg:rgd-final}
    \begin{algorithmic}[1]
        \Require наблюдаемая матрица $Y$, маска $\Omega$, ранг $r$, параметр $\lambda$, начальная матрица $X_0\in\Mcal_r$
        \Ensure приближение $\widehat X$ ранга $r$
        \For{$k=0,1,2,\ldots$}
        \State $G_k \gets \Proj_\Omega(X_k-Y)+\lambda X_k$
        \State $\xi_k \gets \Proj_{T_{X_k}\Mcal_r}(G_k)$
        \State выбрать шаг $\eta_k$
        \State $X_{k+1}\gets \Pi_r(X_k-\eta_k\xi_k)$
        \If{критерий остановки выполнен}
        \State \textbf{break}
        \EndIf
        \EndFor
        \State \Return $X_k$
    \end{algorithmic}
\end{algorithm}

\subsection{Компактная реализация через факторы}

На текущем этапе основной алгоритм записан через матрицу $X$ и ретракцию с помощью усеченного SVD. Это удобно для понимания и для первой реализации, но не идеально для больших матриц: на каждой итерации приходится работать с матрицей размера $m\times n$.

Следующий естественный шаг --- хранить текущую точку в виде
\[
    X=USV^\top,
    \qquad
    U^\top U=I_r,\quad V^\top V=I_r.
\]
Идея состоит в том, чтобы выполнять риманов шаг не как полный переход к большой матрице и обратно, а через факторы $U$, $S$, $V$ и малые матрицы размера порядка $r$. Такой вариант должен быть вычислительно выгоднее, особенно когда $r\ll m,n$.

При этом важно, что это не просто обычный градиентный спуск по независимым переменным $U$, $S$, $V$. Нужно сохранить геометрию многообразия фиксированного ранга: касательное направление, проекцию и возврат на множество матриц ранга $r$. Полная схема такого шага сейчас находится в процессе уточнения. Предполагается, что регуляризация из \eqref{eq:fixed-rank-regularized} переносится в этот вариант так же: в евклидовом градиенте появляется добавка $\lambda X$.

В финальной версии работы компактная реализация будет рассмотрена как основной путь к масштабированию метода. В предварительном тексте мы фиксируем ее идею, но основные эксперименты ниже пока интерпретируем через более прозрачный вариант с SVD-ретракцией.

\newpage

\section{Практическая часть}
\label{sec:empirical}

\subsection{Код и постановка эксперимента}

Для практической части написан отдельный Python-скрипт:
\begin{center}
    \path{matrix_completion_experiments.py}
\end{center}
Скрипт генерирует низкоранговые матрицы, скрывает часть элементов, добавляет шум и сравнивает методы восстановления. Такой синтетический эксперимент нужен как контролируемая проверка: мы знаем настоящую матрицу $X_*$ и можем честно измерять ошибку.

Матрица строится по модели
\[
    X_* = AB^\top,
\]
где $A\in\R^{m\times r}$ и $B\in\R^{n\times r}$. Затем часть элементов отдается алгоритмам для обучения, а часть известных элементов откладывается как test. На наблюдаемых позициях берется
\[
    Y_{ij}=(X_*)_{ij}+\varepsilon_{ij}, \qquad (i,j)\in\Omega,
\]
где $\varepsilon_{ij}$ отвечает за шум.

\subsection{Методы и метрики}

В текущей версии сравниваются:
\begin{enumerate}[label=\arabic*.]
    \item ALS как простой факторизационный baseline;
    \item Soft-Impute как baseline, связанный с ядерной нормой;
    \item риманов градиентный спуск без регуляризации;
    \item риманов градиентный спуск с L2-регуляризацией;
    \item риманов градиентный спуск с выбором $\lambda$ по validation-множеству.
\end{enumerate}

Основные метрики:
\begin{enumerate}[label=\arabic*.]
    \item RMSE и MAE на test-элементах;
    \item относительная ошибка $\fro{\widehat X-X_*}/\fro{X_*}$;
    \item время работы;
    \item число итераций.
\end{enumerate}

Параметр $\lambda$ подбирается двумя способами. Сначала мы смотрим сетку значений, чтобы понять, как регуляризация влияет на качество. Затем используем честный validation-подбор: часть наблюдаемых элементов временно откладывается, $\lambda$ выбирается по ним, а итоговая модель заново обучается на всех наблюдаемых элементах.

\begin{algorithm}[H]
    \caption{Validation-подбор параметра $\lambda$}
    \label{alg:validation-lambda-final}
    \begin{algorithmic}[1]
        \Require наблюдаемые индексы $\Omega_{\mathrm{train}}$, сетка $\Lambda$, доля validation $\rho$
        \Ensure выбранное значение $\lambda_*$ и итоговая матрица $\widehat X$
        \State разбить $\Omega_{\mathrm{train}}$ на $\Omega_{\mathrm{fit}}$ и $\Omega_{\mathrm{val}}$
        \For{$\lambda\in\Lambda$}
        \State обучить RGD на $\Omega_{\mathrm{fit}}$ с параметром $\lambda$
        \State посчитать $\operatorname{RMSE}_{\mathrm{val}}(\lambda)$ на $\Omega_{\mathrm{val}}$
        \EndFor
        \State $\lambda_*\gets \argmin_{\lambda\in\Lambda}\operatorname{RMSE}_{\mathrm{val}}(\lambda)$
        \State заново обучить RGD с $\lambda_*$ на всех индексах $\Omega_{\mathrm{train}}$
        \State оценить итоговую ошибку на test-множестве
        \State \Return $\lambda_*$ и $\widehat X$
    \end{algorithmic}
\end{algorithm}

\subsection{Предварительные результаты}

В текущем запуске использовались синтетические матрицы размера $60\times 60$ истинного ранга $4$. Каждый сценарий запускался в трех повторениях. Менялись уровень шума, доля наблюдаемых элементов, тип инициализации и ранг, который передается алгоритму.

\begin{table}[H]
    \centering
    \small
    \caption{Средние результаты по всем сценариям}
    \label{tab:main-results-final}
    \begin{tabular}{lrrrr}
        \toprule
        Метод                        & RMSE   & Отн. ошибка & Время, c & Итерации \\
        \midrule
        ALS                          & 0.0720 & 0.7307      & 0.1281   & 49.1     \\
        Soft-Impute                  & 0.0420 & 0.4191      & 0.0952   & 110.3    \\
        RGD                          & 0.0443 & 0.4393      & 0.1708   & 195.0    \\
        RGD, $\lambda=10^{-2}$       & 0.0386 & 0.3820      & 0.1733   & 190.1    \\
        RGD, $\lambda=3\cdot10^{-2}$ & 0.0390 & 0.3850      & 0.1520   & 182.8    \\
        RGD, validation $\lambda$    & 0.0359 & 0.3571      & 1.3552   & 1600.4   \\
        \bottomrule
    \end{tabular}
\end{table}

По таблице \ref{tab:main-results-final} видно, что обычный RGD близок к Soft-Impute, но не всегда его обгоняет. После добавления L2-регуляризации качество заметно улучшается. Лучший результат дает validation-подбор $\lambda$, но он намного дороже по времени, потому что для одного сценария приходится обучать несколько моделей.

\begin{table}[H]
    \centering
    \small
    \caption{Средний RMSE при разных уровнях шума}
    \label{tab:noise-results-final}
    \begin{tabular}{lrrrr}
        \toprule
        Шум  & ALS    & Soft-Impute & RGD    & RGD, validation $\lambda$ \\
        \midrule
        0.00 & 0.0450 & 0.0367      & 0.0263 & 0.0240                    \\
        0.02 & 0.0601 & 0.0395      & 0.0368 & 0.0328                    \\
        0.05 & 0.1110 & 0.0498      & 0.0698 & 0.0508                    \\
        \bottomrule
    \end{tabular}
\end{table}

Таблица \ref{tab:noise-results-final} показывает главный вывод по шуму. При нулевом и умеренном шуме риманов подход выглядит очень хорошо. При сильном шуме Soft-Impute остается более устойчивым. Это логично: метод с ядерной нормой сильнее сглаживает сингулярные значения, а fixed-rank метод может переобучаться под шум, если регуляризация подобрана недостаточно удачно.

\begin{table}[H]
    \centering
    \small
    \caption{Сетка L2-регуляризации для RGD}
    \label{tab:l2-sweep-final}
    \begin{tabular}{rrrr}
        \toprule
        $\lambda$       & RMSE   & Отн. ошибка & Время, c \\
        \midrule
        $10^{-4}$       & 0.0442 & 0.4379      & 0.1730   \\
        $3\cdot10^{-4}$ & 0.0439 & 0.4352      & 0.1783   \\
        $10^{-3}$       & 0.0430 & 0.4267      & 0.1802   \\
        $3\cdot10^{-3}$ & 0.0412 & 0.4085      & 0.1807   \\
        $10^{-2}$       & 0.0386 & 0.3820      & 0.1733   \\
        $3\cdot10^{-2}$ & 0.0390 & 0.3850      & 0.1520   \\
        $10^{-1}$       & 0.0463 & 0.4564      & 0.1434   \\
        \bottomrule
    \end{tabular}
\end{table}

Из таблицы \ref{tab:l2-sweep-final} видно, что слишком слабая регуляризация почти не меняет метод, а слишком сильная уже портит качество. В текущем эксперименте лучшие значения находятся около $10^{-2}$ и $3\cdot10^{-2}$.

\begin{table}[H]
    \centering
    \small
    \caption{Как часто validation выбирал разные значения $\lambda$}
    \label{tab:selected-lambda-final}
    \begin{tabular}{rrr}
        \toprule
        $\lambda$       & Число выборов & Доля   \\
        \midrule
        $10^{-4}$       & 15            & 9.3\%  \\
        $3\cdot10^{-4}$ & 4             & 2.5\%  \\
        $10^{-3}$       & 2             & 1.2\%  \\
        $3\cdot10^{-3}$ & 30            & 18.5\% \\
        $10^{-2}$       & 44            & 27.2\% \\
        $3\cdot10^{-2}$ & 42            & 25.9\% \\
        $10^{-1}$       & 25            & 15.4\% \\
        \bottomrule
    \end{tabular}
\end{table}

Validation чаще всего выбирает значения порядка $10^{-2}$ и $3\cdot10^{-2}$, но иногда выбирает и более сильную регуляризацию. Поэтому нельзя просто один раз зафиксировать $\lambda$ и считать, что задача решена: сила регуляризации зависит от шума, доли наблюдений и выбранного ранга.

\subsection{Что это дает для дальнейшей работы}

Предварительные эксперименты поддерживают следующую рабочую гипотезу: риманов градиентный спуск полезен, когда данные действительно близки к низкоранговым и шум не слишком большой. L2-регуляризация заметно улучшает устойчивость, особенно если ранг выбран не идеально.

При этом текущая SVD-ретракция остается вычислительно дорогой для больших матриц. Поэтому следующий шаг --- компактная реализация через $X=USV^\top$. Ожидается, что она будет более оптимальной по времени, потому что вместо полного SVD большой матрицы можно будет работать с факторами и малыми матрицами. Регуляризация при этом сохранится: в градиенте по-прежнему будет добавка $\lambda X$.

После синтетических экспериментов нужно перейти к экономическим или панельным данным. Часть известных значений будет скрываться искусственно, чтобы методы можно было сравнить по восстановлению этих значений. При этом реальные пропуски могут быть неслучайными, поэтому прикладной эксперимент лучше рассматривать как дополнительную проверку, а не как замену синтетике.

\newpage

\section*{Заключение}
\phantomsection
\addcontentsline{toc}{section}{Заключение}

\textbf{Выводы.} В предварительном тексте сформирована основа курсовой работы по геометрическим методам для заполнения матриц. Во введении уточнена исследовательская постановка, в обзоре литературы показан переход от выпуклых методов к fixed-rank и римановым подходам, а в теоретической части добавлены необходимые элементы дифференциальной геометрии: касательное пространство, риманов градиент, проекция и ретракция.

\textbf{Основной результат текущего этапа.} Работа уже имеет не только теоретический, но и практический каркас. Подготовлен код для синтетических экспериментов, где можно отдельно менять шум, долю наблюдений, инициализацию, выбранный ранг и уровень L2-регуляризации в римановом методе. Первые результаты показывают, что регуляризация заметно улучшает качество RGD, а validation-подбор $\lambda$ позволяет честно выбирать параметр без подглядывания в test-множество.

\textbf{Дальнейшие шаги.} На следующем этапе необходимо уточнить и реализовать компактную версию риманова спуска через $X=USV^\top$, подготовить графики сходимости и перейти к проверке на прикладных экономических данных. После этого можно будет формулировать уже не только методические, но и исследовательские выводы о режимах, в которых риманов градиентный спуск оказывается наиболее полезным.

\begin{thebibliography}{99}
    \phantomsection
    \addcontentsline{toc}{section}{Список литературы}

    \bibitem{ivanov-dg}
    А. О. Иванов.
    \newblock Классическая дифференциальная геометрия; Дифференциальная геометрия и топология.
    \newblock Лекционные материалы курса, мехмат МГУ / Teach-in.
    \newblock \href{https://teach-in.ru/course/differential-geometry-and-tensor-analysis-ivanov}{URL1};
    \href{https://teach-in.ru/course/differential-geometry-and-tensor-analysis-ivanov-p2}{URL2}.

    \bibitem{fomenko-dg}
    А. Т. Фоменко.
    \newblock Классическая дифференциальная геометрия.
    \newblock Лекционные материалы курса, мехмат МГУ / Teach-in.
    \newblock \href{https://teach-in.ru/course/classical-difgeom-fomenko/about}{URL}.

    \bibitem{penskoy-seminars}
    А. В. Пенской.
    \newblock Классическая дифференциальная геометрия: семинары.
    \newblock Семинарские материалы курса, мехмат МГУ / Teach-in.
    \newblock \href{https://teach-in.ru/course/kdg-seminars-penskoy/about}{URL}.

    \bibitem{candes-recht}
    Emmanuel J. Cand\`es, Benjamin Recht.
    \newblock Exact Matrix Completion via Convex Optimization.
    \newblock \textit{arXiv preprint} arXiv:0805.4471, 2008.
    \newblock \href{https://arxiv.org/abs/0805.4471}{URL}.

    \bibitem{cai-svt}
    Jian-Feng Cai, Emmanuel J. Cand\`es, Zuowei Shen.
    \newblock A Singular Value Thresholding Algorithm for Matrix Completion.
    \newblock \textit{arXiv preprint} arXiv:0810.3286, 2008.
    \newblock \href{https://arxiv.org/abs/0810.3286}{URL}.

    \bibitem{candes-tao}
    Emmanuel J. Cand\`es, Terence Tao.
    \newblock The Power of Convex Relaxation: Near-Optimal Matrix Completion.
    \newblock \textit{arXiv preprint} arXiv:0903.1476, 2009.
    \newblock \href{https://arxiv.org/abs/0903.1476}{URL}.

    \bibitem{recht}
    Benjamin Recht.
    \newblock A Simpler Approach to Matrix Completion.
    \newblock \textit{Journal of Machine Learning Research}, 12:3413--3430, 2011.
    \newblock \href{https://www.jmlr.org/papers/v12/recht11a.html}{URL}.

    \bibitem{keshavan}
    Raghunandan H. Keshavan, Andrea Montanari, Sewoong Oh.
    \newblock Matrix Completion from a Few Entries.
    \newblock \textit{arXiv preprint} arXiv:0901.3150, 2009.
    \newblock \href{https://arxiv.org/abs/0901.3150}{URL}.

    \bibitem{candes-plan}
    Emmanuel J. Cand\`es, Yaniv Plan.
    \newblock Matrix Completion with Noise.
    \newblock \textit{arXiv preprint} arXiv:0903.3131, 2009.
    \newblock \href{https://arxiv.org/abs/0903.3131}{URL}.

    \bibitem{mazumder}
    Rahul Mazumder, Trevor Hastie, Robert Tibshirani.
    \newblock Spectral Regularization Algorithms for Learning Large Incomplete Matrices.
    \newblock \textit{Journal of Machine Learning Research}, 11:2287--2322, 2010.
    \newblock \href{https://www.jmlr.org/papers/v11/mazumder10a.html}{URL}.

    \bibitem{inductive}
    Prateek Jain, Inderjit S. Dhillon.
    \newblock Provable Inductive Matrix Completion.
    \newblock \textit{arXiv preprint} arXiv:1306.0626, 2013.
    \newblock \href{https://arxiv.org/abs/1306.0626}{URL}.

    \bibitem{overview}
    Mark A. Davenport, Justin Romberg.
    \newblock An Overview of Low-Rank Matrix Recovery from Incomplete Observations.
    \newblock \textit{arXiv preprint} arXiv:1601.06422, 2016.
    \newblock \href{https://arxiv.org/abs/1601.06422}{URL}.

    \bibitem{vandereycken}
    Bart Vandereycken.
    \newblock Low-rank Matrix Completion by Riemannian Optimization --- extended version.
    \newblock \textit{arXiv preprint} arXiv:1209.3834, 2012.
    \newblock \href{https://arxiv.org/abs/1209.3834}{URL}.

    \bibitem{r3mc}
    B. Mishra, R. Sepulchre.
    \newblock R3MC: A Riemannian Three-Factor Algorithm for Low-Rank Matrix Completion.
    \newblock \textit{arXiv preprint} arXiv:1306.2672, 2013.
    \newblock URL: \href{https://arxiv.org/abs/1306.2672}{URL}.

    \bibitem{causal-panel}
    Susan Athey, Mohsen Bayati, Nikolay Doudchenko, Guido W. Imbens, Khashayar Khosravi.
    \newblock Matrix Completion Methods for Causal Panel Data Models.
    \newblock \textit{arXiv preprint} arXiv:1710.10251, 2017.
    \newblock \href{https://arxiv.org/abs/1710.10251}{URL}.

\end{thebibliography}

\end{document}
